\chapter*{ВВЕДЕНИЕ}
\addcontentsline{toc}{chapter}{ВВЕДЕНИЕ}

Рекомендательные системы помогают пользователю выбрать объект из каталога. Перебрать каталог целиком долго, а сформулировать точный запрос трудно. Предметная область работы --- ликёры. Эти напитки сильно различаются по вкусу, крепости и сладости. Покупатель обычно описывает желаемое через знакомые ему образцы вроде <<что-то вроде Baileys, но не яичное>>. Поэтому здесь естественно работает \textbf{контентный} подход, в котором рекомендации строятся по сходству признаков объектов~\cite{lops2011, pazzani2007}.

Цель работы --- разработать рекомендательную систему ликёров. Она должна объединять контентные рекомендации и параметрический поиск. Каждое принятое проектное решение требуется обосновать.

Задачи соответствуют лабораторным работам блока.
\begin{itemize}[leftmargin=1.6\parindent, label*=---]
	\item ЛР1. Сформировать датасет и иерархию предметной области. Дерево должно содержать не менее 25 узлов. Признаки нужны разных типов.
	\item ЛР2. Реализовать и сравнить меры близости для каждого типа признаков. Отдельно нужны мера по дереву и обобщающая мера по всем признакам.
	\item ЛР3. Построить контентную РС в трёх версиях: по затравке, по лайкам и по лайкам с дизлайками. Нужен пользовательский интерфейс.
	\item ЛР4. Реализовать параметрический поиск по фильтрам с проверкой ввода. При пустом результате показывается выдача <<возможно, Вам понравится>>.
	\item РК1. Объединить ЛР3 и ЛР4 в единую систему с общим интерфейсом, фильтрацией рекомендаций и историей запросов.
\end{itemize}

Для каждого этапа в отчёте описано, какие альтернативы рассматривались и почему выбран конкретный вариант. Там, где выбор можно проверить численно, приведены результаты проверки на датасете. Программы проверок входят в репозиторий и воспроизводят все таблицы и рисунки отчёта.
